\section{算法考古：旧论文如何在新硬件上重现光芒}

杨松琳强调“考古”：读旧论文、理解旧机制，然后在新问题和新硬件下重新组合。许多机制早在 2016、2020、2021 年就出现过，但当时缺少合适的硬件、并行算法、开源实现或大规模验证，所以被埋在文献海里。DataNet、Delta Rule、Gated Linear Attention 等都属于这种重新被激活的线索。

\begin{figure}[H]
\centering
\includegraphics[width=0.96\textwidth]{algorithm-archaeology.png}
\caption{算法考古路线：旧论文、旧机制和新硬件重新组合。自制概念图，依据 01:23:12--01:29:50 对谈内容整理。}
\end{figure}

\begin{knowledgebox}{读图：考古不是怀旧，而是寻找可复用工具}
Transformer、Sparse、DeltaNet、Gated Linear、KDA 之间不是直线进步，而是旧机制在新上下文中被重新评估。很多想法不是第一次出现，而是第一次具备 scale up 条件。
\end{knowledgebox}

\subsection{为什么旧工作会被埋没}

本节解释为什么“考古”不是装饰。很多旧论文不是错了，而是在当时缺少让它们变成主流的工程条件。

旧工作被埋没可能有三种原因。第一，配套算法和 kernel 不成熟，别人难以复现或 scale。第二，当时硬件和任务不需要它，价值没有显现。第三，开源代码不好用，社区无法 follow up。杨松琳强调把代码做好，让技术能流传下去，这其实是架构研究影响力的重要组成部分。

\begin{warningbox}{小规模有效不代表大规模有效}
架构变体很多，许多算法在小规模任务上 work，但到大规模 pre-training 会失效。Scaling Ladder、硬件实现和开源复现，是避免小规模幻觉的重要机制。
\end{warningbox}

\subsection{研究品味：知道什么值得做}

杨松琳说，真正大的 challenge 不是技术问题本身，而是不知道要做什么。她的做法是先把领域里值得看的论文都读一遍，形成问题地图，再判断哪个旧机制在今天有价值。比如 DataNet 早在 2021 年就有，但缺硬件效率保证；当任务需求和并行算法成熟后，它就可能重新成为可用组件。

\begin{importantbox}{算法研究的品味}
好的架构研究不是随机组合模块，而是知道历史上有哪些工具、它们为什么没成功、今天哪些约束变了，以及如何用新算法和硬件把旧思想 scale up。
\end{importantbox}

\subsection{开源实现也是研究贡献}

Flash Linear Attention 之所以重要，不只是因为论文，而是因为它让研究者和业界能方便试线性注意力。如果一个想法没有可用实现，后续研究很难 follow；如果实现好用，技术路线更容易流传和被验证。

\begin{knowledgebox}{从论文到路线的三步}
第一步是 idea：数学上或机制上 make sense。第二步是 implementation：能否写成可用、可复现、可优化的代码。第三步是 scaling：能否在更大模型、更真实任务和更硬件友好的条件下成立。
\end{knowledgebox}

\subsection{本章小结}

算法考古提供了一种研究方法：不要只追最新 buzzword，要读旧论文、理解机制、判断哪些想法在今天的硬件和任务下可以重新发光。

